\subsection{Problema de vertex-cover}

Un cubrimiento por vértices (\textit{vertex cover}) de un grafo \((V,E)\) es un subconjunto de vértices \(S\subseteq V\) que cubre todas las aristas, es decir, que toda arista \((u,v) \in E\) tiene al menos un extremo en \(S\): \(u\in S\) o \(v\in S\).

El \concept{problema de vertex-cover} consiste en hallar un cubrimiento por vértices de tamaño mínimo.

\subsubsection{Algoritmo aproximado para vertex-cover}

A pesar de que nadie ha hallado un algoritmo polinomial correcto, una idea bastante intuitiva resulta en un algoritmo de aproximación 2.

La idea es: mientras queden aristas sin cubrir, se elige arbitrariamente una de ellas, se agregan ambos extremos a la solución y se descartan todas las demás aristas incidentes a esos extremos.

\begin{pseudocode}
def aprox_vertex_cover(G: grafo) -> set[vértices]:
    |\(S\)| = |\(\varnothing\)|
    |\(E\)| = |\textrm{copia de las aristas de }\(G\)| 
    while |\(E \neq \varnothing\)|:
        |\(e\)| = |\textrm{eje arbitrario }\((u, v) \in E\)| 
        |\(S\)| = |\(S \cup \{u, v\}\)|
        |\textrm{eliminar cada arista incidente a }\(u\)\textrm{ o }\(v\)\textrm{ de }\(E\)|
    return |\(S\)|
\end{pseudocode}

\subsubsection{Complejidad temporal}



\subsubsection{Razón de aproximación}